
%%%%%%%%%%%%%%%%%%%%%%%%%%%%%%%%%%%%%%%%%%%%%%%%%%%%%%
\chapter*{Disclaimer about the use of GenAI}
%%%%%%%%%%%%%%%%%%%%%%%%%%%%%%%%%%%%%%%%%%%%%%%%%%%%%%

Je déclare que ce mémoire est conforme à la charte d’utilisation de l’IA générative dans le cadre du mémoire à la faculté d’informatique, en accord avec mon ou mes promoteurices.

\subsection*{Usages dans le travail scientifique}

L’ensemble des choix scientifiques, des raisonnements, des développements théoriques et des implémentations relève de mon travail personnel. Les outils d’IA générative n’ont pas été utilisés pour produire directement :
\begin{compactitem}
    \item les questions de recherche ;
    \item les hypothèses ;
    \item la méthodologie expérimentale ;
    \item la conception des expériences ;
    \item l’analyse ou l’interprétation des résultats ;
    \item la modélisation ou la formalisation théorique ;
    \item les implémentations soumises ;
    \item les contributions scientifiques originales présentées dans ce travail.
\end{compactitem}

\subsection*{Usages dans la rédaction du manuscrit}

Les outils d’IA générative n’ont pas été utilisés pour générer de manière substantielle le contenu scientifique, l’argumentation, l’état de l’art, l’analyse des résultats ou les conclusions.

\subsection*{Usages spécifiques de l’IA générative autorisés}

Les promoteurices ont préalablement autorisé l’utilisation d’outils d’IA générative dans le cadre du mémoire pour les usages suivants :
\begin{compactenum}
    \item améliorer la qualité de la langue et corriger la grammaire ainsi que la mise en forme LaTeX du manuscrit (par exemple, via Grammarly ou MS Copilot en mode conversationnel) ;
    \item aider à l'exploration préliminaire de la bibliographie et à l'identification des questions encore peu étudiées ;
    \item aider à la rédaction de scripts Python ou R pour l’analyse et la visualisation des données (par exemple, via MS Copilot en mode conversationnel) ;
    \item aider à la rédaction du code source des implémentations soumises (par exemple, via MS Copilot en mode conversationnel). Attention, il n'est pas autorisé de demander à l'IA générative de produire directement des implémentations complètes, mais seulement d'aider à la rédaction de fragments de code ou de fonctions spécifiques.
\end{compactenum}

\subsection*{Responsabilité}

Je confirme être l’auteur des contributions scientifiques présentées dans ce mémoire et être capable d’expliquer et de défendre l’ensemble de son contenu. Aucune utilisation d’IA générative n’a remplacé mon travail intellectuel.